\scanpage{31}\mspage{VI}
\uline{Demonstration} : Supposons $\Gamma=\partial\mathbb{C}$, avec $\mathbb{C}$ convexe et compact.

On prend $\operatorname{int}\Gamma=\mathring{\mathbb{C}}$ et $\operatorname{ext}\Gamma=\mathbb{R}^{2}-\mathbb{C}$. Il est clair que $\operatorname{int}\Gamma=\mathring{\mathbb{C}}$ est convexe puisque l'interieur d'un ensemble convexe est convexe.

Pour montrer que $\mathbb{R}^{2}-\mathbb{C}$ est connexe on choisit \inserted{le bord $\partial T$ d'}un triangle $T$ (qui vérifie la propiété d'être un J-polygone par des raisons elementaires) et on montre que $(\partial T\cup\operatorname{int}\partial T)-\mathbb{C}$ est connexe. Soit $x\in\partial T\cup\operatorname{int}\partial T$ \deleted{$x\notin\mathbb{C}$,} $(=R_{\partial T})$ \inserted{et $x\notin\mathbb{C}$} et soit $d$ une droite qui passe par $x$. Alors $d\cap\mathbb{C}=\phi$ ou bien $d\cap\mathbb{C}=[z_1,z_2]$\inserted{ avec $z_1,z_2\in\mathbb{C}$}, Dans le premier cas on voit que $x$ peut être \unsure{reliér|relié} à un point de $\partial T$ puisque $d\cap\partial T\ne\phi$.

Si $d\cap\mathbb{C}=[z_1,z_2]$ soit $z_1$ le point \inserted{de $\mathbb{C}$} le plus proche à $x$ sur $d$, alors la demi-droite fermée d'origine \deleted{$\unsure{z|x}$}$z_1$ qui contient $x$ coupe aussi $\partial T$. Donc dans tous le cas $x$ peut être relié a $\partial T$ et cela montre la connexité de $\overset{\inserted{\substack{T\\\Vert}}}{(\partial T\cup\operatorname{int}\partial T)}-\mathbb{C}$, comme $(\mathbb{R}^{2}-T)\cup\partial T$ est connexe \inserted{(fermeture d'un connexe)} \margin{\tikz[x=.08cm,y=-.08cm]{\foreach \x/\y in {1/0,1/1.1,.8/1.5,.75/2.1,.6/2.3,0/2.8,.4/3.4}\fill (\x,\y) circle[radius=.35pt];}}on a montré la propiété \deleted{\unsure{b}}a) de la definition 2. La propiété b) est une application \inserted{simple} du lemme .

\margin{$\bullet$}
L'affirmation qui reste est trivial par definition de convexe, \unsure{\(\blacksquare\)|\deleted{\illegible{1}}} c.q.f.d. .

\margin{$\bullet$}
Si $\Gamma$ est un polygone de $\mathbb{R}^{2}$ on considĕre son enveloppe convexe, $E(\Gamma)$.

\margin{réf \unsure{juste.}}
Il est \uline{clair} que $E(\Gamma)$ vérifie :
\[
\begin{aligned}
\text{a)}\quad & E(\Gamma)=E(S(\Gamma))\quad;\\
\text{b)}\quad & \underset{\inserted{S(\partial(\unsure{E(\Gamma|E\Gamma}))}}{S(\partial(\Gamma))}\subset S(\Gamma)\quad.
\end{aligned}
\]

On veut étudier le lien entre $\Gamma$ et son enveloppe convexe.
% archive sheet number: 30
